% Borg: Extended Abstract for HPG 2026 Poster Track
% 2 pages, SIGGRAPH-style ACM format

\documentclass[sigconf,review]{acmart}

\usepackage{booktabs}
\usepackage{microtype}
\usepackage{xcolor}
\definecolor{borgblue}{HTML}{1A3C5E}

\acmConference[HPG '26]{High-Performance Graphics}{July 17--19, 2026}{Los Angeles, CA}
\acmYear{2026}

\begin{document}

\title{Borg: A Complete Open-Source TBDR GPU Pipeline\\
       with ASIC Tapeout on 5280 Logic Cells}

\author{Andreas Wendleder}
\affiliation{%
  \institution{Independent Researcher}
  \city{Landshut}
  \country{Germany}
}
\email{andreas.wendleder@proton.me}

\begin{abstract}
We present \textsc{Borg}, a Tile-Based Deferred Rendering (TBDR) GPU
implemented in Chisel (a Scala-based hardware description language),
synthesized with an open-source toolchain, and submitted to physical
ASIC fabrication via the Tiny Tapeout programme.
The complete rendering pipeline — vertex shader, triangle clipping,
edge function rasterizer, texture fetch, Z-buffer, and fragment
interpolation — fits within 5{,}280 logic cells of a Lattice iCE40 UP5K FPGA.
To our knowledge, \textsc{Borg} is the first open-source GPU
for which a physical ASIC tapeout has been completed.
We describe the architecture, design decisions, resource budget,
and a path toward hardware-accelerated path tracing.
\end{abstract}

\maketitle

\section{Introduction}

Modern GPU architectures are opaque by design.
Every shipped GPU conceals its microarchitecture, instruction set,
and physical layout behind NDAs or proprietary binaries.
This constrains rendering research in three ways:
results are non-reproducible, architectural experiments require
expensive EDA licenses, and teaching stops at the API boundary.

\textsc{Borg} is a fully open alternative.
Every component — Chisel HDL source, synthesis scripts, place-and-route
configuration, and tapeout submission packages — is publicly available
at \url{https://github.com/gonsolo/Borg} under a permissive licence.

\section{Architecture}

\textsc{Borg} implements the \emph{Tile-Based Deferred Rendering} (TBDR)
paradigm, processing geometry tile-by-tile so that shading is applied
only to visible fragments.
This approach, also used by PowerVR/Imagination Technologies, Apple GPU,
ARM Mali, and Vivante, reduces external memory bandwidth — important
for FPGA designs with limited PSRAM throughput.

The pipeline comprises six hardware stages:
\begin{enumerate}
  \item \textbf{Vertex Shader:} Perspective projection via a hardware
    reciprocal unit (RCP), executing a custom RISC-V-derived 32-bit ISA.
  \item \textbf{Triangle Clipping:} Sutherland-Hodgman clip against
    the tile viewport.
  \item \textbf{Edge Function Unit:} Pineda-style signed area test,
    one pixel per clock.
  \item \textbf{Texture Fetch:} Nearest-neighbour UV sampling from PSRAM
    via a shared memory controller.
  \item \textbf{Fragment Interpolation + Z-Buffer:} Perspective-correct
    barycentric attribute interpolation with depth test.
  \item \textbf{Tile Buffer:} On-chip BRAM with Morton (Z-curve) encoding
    for cache-friendly 2-D access.
\end{enumerate}

A command FIFO decouples the CPU from the rendering pipeline.
A SystemRDL register map (generated via PeakRDL-chisel) provides
a clean software interface.

\section{Resource Budget}

\begin{table}[h]
  \centering
  \caption{iCE40 UP5K resource utilisation (pico-ice board).}
  \begin{tabular}{@{}lrr@{}}
    \toprule
    Resource & Used & Available \\
    \midrule
    Logic Cells & $\sim$5\,100 & 5\,280 \\
    4\,kb BRAMs & \todo{N} & 30 \\
    DSPs        & 0 & 8 \\
    \bottomrule
  \end{tabular}
\end{table}

The iCE40 UP5K provides no DSP blocks usable for FP multiply-add;
all floating-point computation is implemented in soft logic.

\section{ASIC Tapeouts}

\begin{table}[h]
  \centering
  \caption{Physical tapeouts submitted as part of the Borg project.}
  \begin{tabular}{@{}llll@{}}
    \toprule
    Run & Process & Content & Submitted \\
    \midrule
    GF180    & GF 180\,nm   & Diffuse shader ASIC    & 2022 \\
    TTIHP26a & IHP 130\,nm  & FP adder peripheral    & Mar 2026 \\
    TTSKY26a & SKY130       & Triangle rasteriser    & May 2026 \\
    \bottomrule
  \end{tabular}
\end{table}

The GF180 tapeout was funded by the Google Open-Source Silicon Programme;
TTIHP26a and TTSKY26a via the Tiny Tapeout programme.
Chips for TTIHP26a are expected in July 2026.

\section{Comparison with Related Open GPU Projects}

\textsc{Nyuzi}~\cite{Bush2015Nyuzi} is an open GPGPU targeting compute
workloads (no rasterisation pipeline, no tapeout).
\textsc{Vortex}~\cite{Vortex2021} is a RISC-V GPGPU targeting research
on GPU computing; it targets Xilinx FPGAs and has not been taped out.
\textsc{Borg} differs in focus (graphics, not compute),
architecture (TBDR, not SIMD lanes), and,
critically, in having submitted physical silicon.

\section{Current Status and Future Work}

Phase 1 (Steps 1–20: the full rendering pipeline) is complete.
Phase 2 (GPU autonomy — DMA engine, write port, vertex sequencer)
is in progress and will remove CPU orchestration from the render loop.
Phase 3 plans a Linux-capable CPU integration, and Phase 4 targets
a minimal Mesa Vulkan driver.

The longer-term goal is integration with \textsc{Gonzales}~\cite{Wendleder2024Gonzales},
a production path tracer, to create the first fully open-source hardware
path tracing system.

\section{Acknowledgements}

The Borg project is funded by the NGI0 Commons Fund established by
NLnet Foundation with financial support from the European Commission's
Next Generation Internet programme
(Grant Agreement No.~101135429, DG CONNECT),
and by the Swiss State Secretariat for Education, Research and
Innovation (SERI).

\bibliographystyle{ACM-Reference-Format}
\bibliography{../bibliography}

\end{document}
